% GENERATED FILE. Edit canonical lessons, then run npm run generate:books.
% canonical-source-hash: fnv1a64:8cebd174cedef279
% canonical-lessons: BN-C285-shasano, BN-C285-ghush, BN-C285-phasano, BN-C285-lut-kora, BN-C285-protarona

\chapter{Threats and Tricks}
\label{ch:bn-threats-and-tricks}
\hlchaptermodality{285}

\begin{chapteropening}
\textbf{By the end of this chapter:} \emph{I can talk about threats, bribes, frame-ups and looting.}

Complete the last lesson of chapter 285: I can talk about threats, bribes, frame-ups and looting.
\end{chapteropening}

\section[\texorpdfstring{\bn{শাসানো} (shāsāno)}{শাসানো (shāsāno)}]{\bn{শাসানো} (shāsāno) — to threaten; to warn sternly}

\label{lesson:BN-C285-shasano}



\begin{quote}
Before the new one: say the Bengali for a prisoner, then the Bengali for bail.
\end{quote}

\subsection*{You'll want to know: \bn{শাসানো}}
\textbf{\bn{শাসানো}} — \emph{shāsāno} — \textquotedblleft{}to threaten; to warn sternly\textquotedblright{}.

\begin{quote}
\bn{ডাকাতরা} \bn{দোকানদারকে} \bn{শাসাল}, \bn{কিন্তু} \bn{তিনি} \bn{ভয়} \bn{পেলেন} \bn{না।}
\end{quote}

\begin{quote}
\emph{ḍākātrā dokāndārke shāsālô, kintu tini bhôy pelen nā.}
\end{quote}

\begin{quote}
The robbers threatened the shopkeeper, but he wasn't afraid.
\end{quote}

Past (he/she …): \bn{সেই} \bn{ডাকাত} \bn{দোকানদারকে} \bn{শাসাল}, \emph{sei ḍākāt dokāndārke shāsālô}.

\begin{cousinweb}
\textbf{\bn{শাসানো}} (\emph{shāsāno}) comes down from Sanskrit \emph{śāsati}, 'instructs, punishes, rules' (CDIAL 12419, which lists Bengali \emph{sāsā}, 'to punish', and notes the -ānō verb, 'to threaten'). Its root \emph{śās-} is also the root of \textbf{\bn{শাস্তি}}, a punishment: to threaten is to hold the punishment over someone.
\end{cousinweb}

\begin{grammarlens}[title={What you've built}]
One more word for telling what happened.
\end{grammarlens}

\subsection*{Your turn}
\begin{itemize}
\raggedright
  \item \emph{Say it:} \emph{shāsāno}
  \item \emph{Say it:} \emph{shāsāno}, once more
  \item \emph{Say it:} \emph{ḍākātrā dokāndārke shāsālô, kintu tini bhôy pelen nā.} — The robbers threatened the shopkeeper, but he wasn't afraid.
  \item \emph{Recall:} say the Bengali for a prisoner, then the Bengali for bail, then say \emph{shāsāno} again
\end{itemize}

\begin{tcolorbox}[breakable,colback=teal!4,colframe=teal!35!black,title={Before you move on}]
What does \bn{শাসানো} mean? (\textquotedblleft{}To threaten; to warn sternly\textquotedblright{}.) Say it once more.
\end{tcolorbox}
\section[\texorpdfstring{\bn{ঘুষ} (ghush)}{ঘুষ (ghush)}]{\bn{ঘুষ} (ghush) — a bribe}

\label{lesson:BN-C285-ghush}



\begin{quote}
Before the new one: say the Bengali for bail, then the Bengali for to threaten.
\end{quote}

\subsection*{You'll want to know: \bn{ঘুষ}}
\textbf{\bn{ঘুষ}} — \emph{ghush} — \textquotedblleft{}a bribe\textquotedblright{}.

\begin{quote}
\bn{চোর} \bn{দারোগাকে} \bn{টাকা} \bn{দিতে} \bn{চাইল}, \bn{কিন্তু} \bn{তিনি} \bn{ঘুষ} \bn{নিলেন} \bn{না।}
\end{quote}

\begin{quote}
\emph{chor dārogāke ṭākā dite chāilô, kintu tini ghush nilen nā.}
\end{quote}

\begin{quote}
The thief tried to give the inspector money, but he would not take a bribe.
\end{quote}

\begin{cousinweb}
\textbf{\bn{ঘুষ}} (\emph{ghush}) is filed in CDIAL 4499 under a Middle Indo-Aryan \emph{ghūssā}, 'a present', which lists Bengali \emph{ghus}, 'a bribe': a present with a purpose.
\end{cousinweb}

\begin{grammarlens}[title={What you've built}]
One more word for telling what happened.
\end{grammarlens}

\subsection*{Your turn}
\begin{itemize}
\raggedright
  \item \emph{Say it:} \emph{ghush}
  \item \emph{Say it:} \emph{ghush}, once more
  \item \emph{Say it:} \emph{chor dārogāke ṭākā dite chāilô, kintu tini ghush nilen nā.} — The thief tried to give the inspector money, but he would not take a bribe.
  \item \emph{Recall:} say the Bengali for bail, then the Bengali for to threaten, then say \emph{ghush} again
\end{itemize}

\begin{tcolorbox}[breakable,colback=teal!4,colframe=teal!35!black,title={Before you move on}]
What does \bn{ঘুষ} mean? (\textquotedblleft{}A bribe\textquotedblright{}.) Say it once more.
\end{tcolorbox}
\section[\texorpdfstring{\bn{ফাঁসানো} (phā̃sāno)}{ফাঁসানো (phā̃sāno)}]{\bn{ফাঁসানো} (phā̃sāno) — to frame; to trap (someone)}

\label{lesson:BN-C285-phasano}



\begin{quote}
Before the new one: say the Bengali for to threaten, then the Bengali for a bribe.
\end{quote}

\subsection*{You'll want to know: \bn{ফাঁসানো}}
\textbf{\bn{ফাঁসানো}} — \emph{phā̃sāno} — \textquotedblleft{}to frame; to trap (someone)\textquotedblright{}.

\begin{quote}
\bn{আসল} \bn{চোর} \bn{আমার} \bn{ভাইকে} \bn{ফাঁসাতে} \bn{চাইল।}
\end{quote}

\begin{quote}
\emph{āsôl chor āmār bhāike phā̃sāte chāilô.}
\end{quote}

\begin{quote}
The real thief tried to frame my brother.
\end{quote}

Past (he/she …): \bn{সেই} \bn{চোর} \bn{তাকে} \bn{ফাঁসাল}, \emph{sei chor tāke phā̃sālô}.

\begin{cousinweb}
\textbf{\bn{ফাঁসানো}} (\emph{phā̃sāno}) comes down from Sanskrit \emph{spāśayati}, 'causes to be bound' (CDIAL 13814, which lists Bengali \emph{phā̃sā}, 'to be ensnared'). A \bn{ফাঁস} is a noose: to frame someone is to slip a noose on them.
\end{cousinweb}

\begin{grammarlens}[title={What you've built}]
One more word for telling what happened.
\end{grammarlens}

\subsection*{Your turn}
\begin{itemize}
\raggedright
  \item \emph{Say it:} \emph{phā̃sāno}
  \item \emph{Say it:} \emph{phā̃sāno}, once more
  \item \emph{Say it:} \emph{āsôl chor āmār bhāike phā̃sāte chāilô.} — The real thief tried to frame my brother.
  \item \emph{Recall:} say the Bengali for to threaten, then the Bengali for a bribe, then say \emph{phā̃sāno} again
\end{itemize}

\begin{tcolorbox}[breakable,colback=teal!4,colframe=teal!35!black,title={Before you move on}]
What does \bn{ফাঁসানো} mean? (\textquotedblleft{}To frame; to trap (someone)\textquotedblright{}.) Say it once more.
\end{tcolorbox}
\section[\texorpdfstring{\bn{লুট} \bn{করা} (luṭ kôrā)}{লুট করা (luṭ kôrā)}]{\bn{লুট} \bn{করা} (luṭ kôrā) — to loot; to plunder}

\label{lesson:BN-C285-lut-kora}



\begin{quote}
Before the new one: say the Bengali for a bribe, then the Bengali for to frame.
\end{quote}

\subsection*{You'll want to know: \bn{লুট} \bn{করা}}
\textbf{\bn{লুট} \bn{করা}} — \emph{luṭ kôrā} — \textquotedblleft{}to loot; to plunder\textquotedblright{}.

\begin{quote}
\bn{ডাকাতরা} \bn{রাতে} \bn{বাজার} \bn{লুট} \bn{করল।}
\end{quote}

\begin{quote}
\emph{ḍākātrā rāte bājār luṭ kôrlô.}
\end{quote}

\begin{quote}
The robbers looted the market at night.
\end{quote}

Past (he/she …): \bn{সেই} \bn{ডাকাত} \bn{দোকান} \bn{লুট} \bn{করল}, \emph{sei ḍākāt dokān luṭ kôrlô}.

\begin{cousinweb}
\textbf{\bn{লুট} \bn{করা}} (\emph{luṭ kôrā}) puts \bn{লুট}, 'plunder' — CDIAL 11078, a Middle Indo-Aryan \emph{luṭṭati}, which lists Bengali \emph{luṭā}, 'plunders' — before \bn{করা}. English \emph{loot} was borrowed from this Indian word.
\end{cousinweb}

\begin{grammarlens}[title={What you've built}]
One more word for telling what happened.
\end{grammarlens}

\subsection*{Your turn}
\begin{itemize}
\raggedright
  \item \emph{Say it:} \emph{luṭ kôrā}
  \item \emph{Say it:} \emph{luṭ kôrā}, once more
  \item \emph{Say it:} \emph{ḍākātrā rāte bājār luṭ kôrlô.} — The robbers looted the market at night.
  \item \emph{Recall:} say the Bengali for a bribe, then the Bengali for to frame, then say \emph{luṭ kôrā} again
\end{itemize}

\begin{tcolorbox}[breakable,colback=teal!4,colframe=teal!35!black,title={Before you move on}]
What does \bn{লুট} \bn{করা} mean? (\textquotedblleft{}To loot; to plunder\textquotedblright{}.) Say it once more.
\end{tcolorbox}
\section[\texorpdfstring{\bn{প্রতারণা} (prôtārônā)}{প্রতারণা (prôtārônā)}]{\bn{প্রতারণা} (prôtārônā) — fraud; deception}

\label{lesson:BN-C285-protarona}



\begin{quote}
Before the new one: say the Bengali for to frame, then the Bengali for to loot.
\end{quote}

\subsection*{You'll want to know: \bn{প্রতারণা}}
\textbf{\bn{প্রতারণা}} — \emph{prôtārônā} — \textquotedblleft{}fraud; deception\textquotedblright{}.

\begin{quote}
\bn{একটা} \bn{লোক} \bn{প্রতারণা} \bn{করে} \bn{গ্রামের} \bn{মানুষের} \bn{টাকা} \bn{নিল।}
\end{quote}

\begin{quote}
\emph{ekṭā lok prôtārônā kôre grāmer mānusher ṭākā nilô.}
\end{quote}

\begin{quote}
A man tricked the villagers out of their money.
\end{quote}

\begin{cousinweb}
\textbf{\bn{প্রতারণা}} (\emph{prôtārônā}) is Sanskrit \emph{pratāraṇa}, from the causative of \emph{pra-√tṝ}, which Monier-Williams gives both as 'ferrying across' and as 'deceiving': taking someone across, and taking them in.
\end{cousinweb}

\begin{grammarlens}[title={What you've built}]
One more word for telling what happened.
\end{grammarlens}

\subsection*{Your turn}
\begin{itemize}
\raggedright
  \item \emph{Say it:} \emph{prôtārônā}
  \item \emph{Say it:} \emph{prôtārônā}, once more
  \item \emph{Say it:} \emph{ekṭā lok prôtārônā kôre grāmer mānusher ṭākā nilô.} — A man tricked the villagers out of their money.
  \item \emph{Recall:} say the Bengali for to frame, then the Bengali for to loot, then say \emph{prôtārônā} again
\end{itemize}

\begin{tcolorbox}[breakable,colback=teal!4,colframe=teal!35!black,title={Before you move on}]
What does \bn{প্রতারণা} mean? (\textquotedblleft{}Fraud; deception\textquotedblright{}.) Say it once more.
\end{tcolorbox}
